%%% FIELD concentration-statement
Fix \(Q\in\mathcal P_{\mathrm{att}}\), condition on the projected pilots, put \(x_\eta=\log(8/\eta)\) and \(M_h=mh^d\), and let \(T_*^\circ=\mathbb E_Q(T_*\mid\bar\pi,\bar\mu)\).  There are constants \(C_0,C_1<\infty\), depending only on \((d,\kappa,\alpha,\beta,\gamma,L_\pi,L_\mu)\) and the fixed basis normalizations, such that, for every finite integer \(k\ge1\), whenever \(M_h\ge C_0x_\eta\),
\[
 \Pr_Q\!\left[
 \lVert T_*-T_*^\circ\rVert_2>C_1\left\{
 \sqrt{\frac{x_\eta}{M_h}\left(1+\frac{k}{M_h}\right)}
 +\frac{x_\eta}{M_h}
 +\frac{\sqrt{k}\,x_\eta^{3/2}}{M_h^{3/2}}
 +\frac{kx_\eta^2}{M_h^2}
 \right\}\,\middle|\,\bar\pi,\bar\mu\right]
 \le\eta/2.
\]
The bound is uniform in the realized pilots and in \(J_\pi,J_\mu\), and it imposes no rank condition relating \(k\) to \(M_h\).

%%% FIELD concentration-proof
Write \(r_p=p_\gamma^2+p_\gamma\), the fixed number of scalar coordinates in the concatenated main-block statistic, and fix one coordinate \(j\).  The clipping radii in Definition \(\mathrm{def{:}private\text{-}higher\text{-}order\text{-}learner}\) are the suprema of the corresponding kernel norms over the projected pilot ranges, so the displayed clipping maps leave \(L\) and \(H\) unchanged.  Assumption \(\mathrm{ass{:}bounded\text{-}outcomes}\), the range projection in that definition, Assumption \(\mathrm{ass{:}overlap}\), and the fixed-dimensional bound on \(\rho_h\) give
\[
 |L_j(O)|\le Ch^{-d}.
\]
Assumption \(\mathrm{ass{:}uniform\text{-}design}\) and \(\operatorname{Leb}(C_h)=h^d\) then give
\[
 \mathbb E_Q L_j(O)^2\le
 C h^{-2d}\Pr_Q(X\in C_h)=C h^{-d}.                         \tag{6}
\]
All constants in this proof are uniform in the realized projected pilots because those pilots lie in their deterministic range sets.

If \(Z\) is centered, \(|Z|\le a\), and \(0\le t<3/a\), expansion of the exponential series, the inequality \(r!\ge 2\,3^{r-2}\) for \(r\ge2\), and \(\mathbb E|Z|^r\le a^{r-2}\mathbb EZ^2\) imply
\[
 \mathbb E e^{tZ}
 \le 1+\frac{t^2\mathbb EZ^2}{2}
       \sum_{r=2}^{\infty}(ta/3)^{r-2}
 \le \exp\!\left\{\frac{t^2\mathbb EZ^2}{2(1-ta/3)}\right\}. \tag{7}
\]
Applying (7) to the independent centered linear terms, multiplying their moment-generating functions, and optimizing the resulting Bernstein bound gives, for every \(p\ge2\),
\[
 \left\lVert m^{-1}\sum_{i=1}^m
       \{L_j(O_i)-\mathbb E_Q L_j(O_i)\}\right\rVert_{L^p}
 \le C\left\{\sqrt{\frac p{M_h}}+\frac p{M_h}\right\}.       \tag{8}
\]

We next handle the ordered quadratic statistic.  Symmetrize its scalar kernel and take the Hoeffding decomposition into two first-order projections and a canonical kernel \(H_{2,j}\).  The partial projection kernels associated with the boundary-corrected spline-wavelet basis are uniformly bounded on \(L^\infty(C_h)\).  Indeed, resolution ordering writes a partial projection as all complete resolution levels plus a subset of one final level.  The complete levels form the boundary-corrected scaling projection, whose integral kernel has uniformly bounded absolute integral; at the partial final level, compact support and finite overlap give the same bound term by term.  After rescaling the cube, this proves
\[
 \sup_{k\ge1}\sup_{x\in C_h}
 \int_{C_h}\left|b_{h,k}(x)^{\mathsf T}b_{h,k}(z)\right|\,dz
 \le C.                                                        \tag{9}
\]
Consequently each first-order projection has the same envelope and second-moment bounds as (6), and (8) applies to its empirical average.

For the canonical component, orthonormality, uniform design, the deterministic multiplier bounds already used above, and
\[
 \iint_{C_h^2}\{b_{h,k}(x)^{\mathsf T}b_{h,k}(z)\}^2\,dx\,dz=k
\]
give
\[
 \begin{aligned}
 \mathbb E_Q H_{2,j}(O,O')^2&\le Ckh^{-2d},\\
 \lVert H_{2,j}\rVert_{L^2(P_Q)\to L^2(P_Q)}&\le Ch^{-d},\\
 \sup_o\mathbb E_Q\{H_{2,j}(o,O')^2+H_{2,j}(O',o)^2\}&\le Ckh^{-3d},\\
 \lVert H_{2,j}\rVert_\infty&\le Ckh^{-2d}.
 \end{aligned}                                                \tag{10}
\]
The second line is the operator norm of the integral operator induced by the canonical kernel.  These four estimates hold for every finite \(k\); none compares \(k\) with \(M_h\).

For clarity, we derive the needed canonical moment estimate directly.  For an even integer \(2r\), expand the \(2r\)-th moment of \(\sum_{i\ne j}H_{2,j}(O_i,O_j)\) and associate each monomial with its multigraph on the sample labels.  Canonicality makes every term having a degree-one vertex vanish.  Repeated Cauchy--Schwarz removes, respectively, doubled edges, cycles, vertices of degree at least three, and fully repeated edges; the four successive numerical inputs are the four lines of (10).  The corresponding label counts are bounded by \((Cr)^r\), \((Cr)^{2r}\), \((Cr)^{3r}\), and \((Cr)^{4r}\).  Taking \(2r\)-th roots, and then replacing an arbitrary \(p\ge2\) by the next even integer, yields
\[
 \begin{aligned}
 \left\lVert\sum_{i\ne j}H_{2,j}(O_i,O_j)\right\rVert_{L^p}
 \le C\{&m\sqrt{kp}\,h^{-d}+mp h^{-d}\\
          &+\sqrt{mk h^{-3d}}\,p^{3/2}+kh^{-2d}p^2\}.
 \end{aligned}                                                \tag{11}
\]
This is a finite moment expansion, so it requires neither an asymptotic approximation nor a restriction on finite \(k\).  Since \(n\ge12\) gives \(m=\lfloor n/3\rfloor\ge4\), division by \(m(m-1)\), followed by \(M_h=mh^d\), gives for the normalized canonical term \(U_{2,j}\)
\[
 \lVert U_{2,j}\rVert_{L^p}
 \le C\left\{
 \frac{\sqrt{kp}}{M_h}+\frac p{M_h}
 +\frac{\sqrt{k}\,p^{3/2}}{M_h^{3/2}}
 +\frac{kp^2}{M_h^2}\right\}.                               \tag{12}
\]
Combining (8), its two first-order analogues, and (12) bounds one coordinate of \(T_*-T_*^\circ\) by
\[
 C\left\{
 \sqrt{\frac p{M_h}\left(1+\frac{k}{M_h}\right)}
 +\frac p{M_h}
 +\frac{\sqrt{k}\,p^{3/2}}{M_h^{3/2}}
 +\frac{kp^2}{M_h^2}\right\}.                               \tag{13}
\]
Because \(x_\eta=\log(8/\eta)>\log 32>2\), choose \(p=c_{r_p}x_\eta\), where the fixed constant \(c_{r_p}\) is large enough for a union bound over the \(r_p\) coordinates.  Markov's inequality in the form
\[
 \Pr\{|Z|>e\lVert Z\rVert_{L^p}\}\le e^{-p}
\]
then makes the conditional failure probability at most \(\eta/2\).  Equivalence of Euclidean and coordinatewise norms in the fixed output dimension changes only \(C_1\).  Taking \(C_0\) large enough for the displayed effective-sample premise completes the proof.  Inspection of (6)--(13) shows that the only condition on the basis dimension was \(1\le k<\infty\).

%%% FIELD remainder-statement
There are constants \(C_B,C_S,C_G\), depending only on \((d,\kappa,\alpha,\beta,\gamma,L_\pi,L_\mu,L_\tau)\) and the fixed basis normalizations, such that for every \(Q\in\mathcal P_{\mathrm{att}}\), every \(0<h\le r_0\), and every finite integer \(k\ge1\) with \(m h^d\ge C_S^2\log(8/\eta)\), the projected private pilots satisfy the same H\"older-radius and range restrictions as the primitive nuisances and
\[
 \Pr_Q\!\left\{\left|\widehat\tau_{\mathrm{DP}}(x_0;h,k)-\tau_Q(x_0)\right|>
 B(h,k)+S(h,k,\eta)+G(h,k,\eta,\epsilon_n,\delta_n)\right\}\le\eta.
\]
The three displayed summands respectively bound localization plus higher-order nuisance bias, sampling fluctuation of the linear and degenerate terms, and all propagated Gaussian release noise.  No rank condition relating \(k\) to \(m h^d\) is required.

%%% FIELD remainder-proof
Fix \(Q\in\mathcal P_{\mathrm{att}}\), \(0<h\le r_0\), a finite integer \(k\ge1\), and the pilot dimensions, and condition on the two pilot releases.  By the metric-projection step in Definition \(\mathrm{def{:}private\text{-}higher\text{-}order\text{-}learner}\), the projected pilots lie in the same H\"older balls and ranges as the true nuisances.  By \(\mathrm{lem{:}localized\text{-}spline\text{-}population\text{-}algebra}\), with \(\ell=h/k^{1/d}\), there is \(\vartheta_\tau\) such that
\[
 \rho_h(x_0)^{\mathsf T}\vartheta_\tau=\tau_Q(x_0),\quad
 \lVert R^\circ-Q^\circ\vartheta_\tau\rVert_2\le Cb_0,
 \quad
 \lVert Q^\circ-Q^0\rVert_{\mathrm{op}}\le C\ell^{2\alpha}, \tag{14}
\]
where
\[
 b_0=h^\gamma+\ell^{\alpha+\beta}+h^\gamma\ell^\alpha,\quad
 \lambda_{\min}(Q^0)\ge\lambda_0:=\kappa(1-\kappa),\quad
 \lVert R^\circ\rVert_2\le C.                                \tag{15}
\]
The positivity \(\lambda_0>0\) follows from Assumption \(\mathrm{ass{:}overlap}\), since \(0<\kappa<1/2\).  Also \(\ell\le h\le r_0\) because \(k\ge1\).

Put \(x_\eta=\log(8/\eta)\), \(M_h=mh^d\), and
\[
 a_0=\sqrt{\frac{x_\eta}{M_h}\left(1+\frac{k}{M_h}\right)}
 +\frac{x_\eta}{M_h}
 +\frac{\sqrt{k}\,x_\eta^{3/2}}{M_h^{3/2}}
 +\frac{kx_\eta^2}{M_h^2},
\]
\[
 g_0=\frac{\sqrt{\log(1/\delta_n)x_\eta}}{m\epsilon_n}
       \{h^{-d}+kh^{-2d}\}.                                  \tag{16}
\]
Choose \(C_S\ge\sqrt{C_0}\), where \(C_0\) is from \(\mathrm{lem{:}localized\text{-}score\text{-}concentration\text{-}unrestricted\text{-}rank}\).  The hypothesis \(M_h\ge C_S^2x_\eta\) permits application of that lemma for the present arbitrary finite \(k\).  Hence, conditionally on the pilots, there is an event \(E_S\) with probability at least \(1-\eta/2\) on which
\[
 \lVert T_*-T_*^\circ\rVert_2\le Ca_0.                       \tag{17}
\]

The released main-block Gaussian vector has fixed dimension \(p_\gamma^2+p_\gamma\).  The Gaussian moment-generating identity and a coordinate union bound give
\[
 \Pr\{\lVert Z_*\rVert_2>C\sqrt{x_\eta}\}\le\eta/2.        \tag{18}
\]
The declared ranges \(0<\delta_n<1/2\) and \(0<\epsilon_n\le1\) imply
\[
 c_{\delta_n}\le C\sqrt{\log(1/\delta_n)},\quad
 (\epsilon_n^\circ)^{-1}\le2\epsilon_n^{-1}.                 \tag{19}
\]
The sensitivity conclusion of \(\mathrm{lem{:}private\text{-}release\text{-}and\text{-}sensitivity}\), (18), and (19) therefore give an event \(E_G\), conditionally independent of \(E_S\), such that
\[
 \Pr(E_G\mid\bar\pi,\bar\mu)\ge1-\eta/2,\quad
 \lVert\widetilde T_*-T_*\rVert_2\le Cg_0.                  \tag{20}
\]
On \(E=E_S\cap E_G\), a union bound gives conditional probability at least \(1-\eta\), and (17)--(20) give
\[
 \lVert\widetilde Q-Q^\circ\rVert_{\mathrm{op}}
 +\lVert\widetilde R-R^\circ\rVert_2\le e_0:=C(a_0+g_0).     \tag{21}
\]

We now separate the stable and large-error cases.  Suppose first that
\[
 C\ell^{2\alpha}+e_0\le\lambda_0/2.                           \tag{22}
\]
For every unit vector \(v\), the triangle inequality and (14), (21) give
\[
 \lVert\widetilde Qv\rVert_2
 \ge \lVert Q^0v\rVert_2
      -\lVert(\widetilde Q-Q^\circ)v\rVert_2
      -\lVert(Q^\circ-Q^0)v\rVert_2.
\]
Taking the infimum over such \(v\) and using (22) yields
\[
 \sigma_{\min}(\widetilde Q)
 \ge\lambda_{\min}(Q^0)-\lVert Q^\circ-Q^0\rVert_{\mathrm{op}}
     -\lVert\widetilde Q-Q^\circ\rVert_{\mathrm{op}}
 \ge\lambda_0/2>\lambda_0/4.                                \tag{23}
\]
Thus the spectral floor is inactive, \(\widehat Q_{\mathrm{DP}}=\widetilde Q\), and inversion is valid with
\(\lVert\widetilde Q^{-1}\rVert_{\mathrm{op}}\le2/\lambda_0\).  Using (14), (15), (21), and the fixed bound on \(\rho_h(x_0)\),
\[
 \begin{aligned}
 \left|\widehat\tau_{\mathrm{DP}}(x_0;h,k)-\tau_Q(x_0)\right|
 &\le \lVert\rho_h(x_0)\rVert_2\lVert\widetilde Q^{-1}\rVert_{\mathrm{op}}
       \lVert\widetilde R-\widetilde Q\vartheta_\tau\rVert_2\\
 &\le C\{b_0+a_0+g_0\}.
 \end{aligned}                                                \tag{24}
\]

If (22) fails, then either \(C\ell^{2\alpha}\ge\lambda_0/4\) or \(e_0\ge\lambda_0/4\).  In the first case, since \(\alpha,\beta>0\), raising the resulting positive lower bound for \(\ell\) to the power \(\alpha+\beta\) gives \(\ell^{\alpha+\beta}\ge c>0\).  In the second case, \(a_0+g_0\ge c>0\).  Hence in either case
\[
 b_0+a_0+g_0\ge c_*>0.                                      \tag{25}
\]
The singular-value floor always gives
\(\lVert\widehat Q_{\mathrm{DP}}^{-1}\rVert_{\mathrm{op}}\le4/\lambda_0\).  Equations (15) and (21), together with \(|\tau_Q(x_0)|\le1\), therefore imply on \(E\)
\[
 \left|\widehat\tau_{\mathrm{DP}}(x_0;h,k)-\tau_Q(x_0)\right|
 \le C\{1+a_0+g_0\}
 \le C'\{b_0+a_0+g_0\},                                    \tag{26}
\]
where the last inequality uses (25).  Choose \(C_B,C_S,C_G\) at least as large as the constants in (24) and (26), retaining \(C_S\ge\sqrt{C_0}\).  Definitions \(\mathrm{def{:}bias\text{-}envelope}\), \(\mathrm{def{:}sampling\text{-}envelope}\), and \(\mathrm{def{:}privacy\text{-}envelope}\) then identify
\[
 B(h,k)=C_Bb_0,\quad S(h,k,\eta)=C_Sa_0,\quad
 G(h,k,\eta,\epsilon_n,\delta_n)=C_Gg_0.
\]
Thus the asserted error inequality holds on \(E\).  Its conditional failure probability is at most \(\eta\), uniformly in the pilot realization and for every finite \(k\).  Integrating over the two pilot blocks proves the unconditional bound, including all pilot and main-block Gaussian randomness.

%%% FIELD honest-statement
For every \(0<h\le r_0\), every finite integer \(k\ge1\) satisfying \(m h^d\ge C_S^2\log(8/\eta)\), and every \(J_\pi,J_\mu\ge1\), the released set \(C_n(h,k)\) is replacement \((\epsilon_n,\delta_n)\)-differentially private and satisfies
\[
 \inf_{Q\in\mathcal P_{\mathrm{att}}}
 \Pr_Q\{\tau_Q^{\mathrm c}(x_0)\in C_n(h,k)\}\ge1-\eta.
\]
Coverage probability integrates the sampling law, all pilot and moment Gaussian noises, and the deterministic sample split.  No rank condition relating \(k\) to \(m h^d\) is required.

%%% FIELD honest-proof
Fix \(0<h\le r_0\), a finite integer \(k\ge1\) satisfying the effective-sample condition, and \(J_\pi,J_\mu\ge1\).  By \(\mathrm{lem{:}private\text{-}release\text{-}and\text{-}sensitivity}\), the learner is replacement \((\epsilon_n,\delta_n)\)-differentially private; Definition \(\mathrm{def{:}private\text{-}confidence\text{-}interval}\) makes interval formation post-processing, so \(C_n(h,k)\) has the same privacy guarantee.

For \(Q\in\mathcal P_{\mathrm{att}}\), define
\[
 E_Q=\left\{
 \left|\widehat\tau_{\mathrm{DP}}(x_0;h,k)-\tau_Q(x_0)\right|
 \le B(h,k)+S(h,k,\eta)+G(h,k,\eta,\epsilon_n,\delta_n)
 \right\}.
\]
The proved \(\mathrm{lem{:}uniform\text{-}remainder\text{-}unrestricted\text{-}rank}\) gives \(\Pr_Q(E_Q)\ge1-\eta\).  Proposition \(\mathrm{prop{:}attainable\text{-}inclusion}\) places \(Q\) in \(\mathcal P_{\mathrm{full}}\), and \(\mathrm{lem{:}pointwise\text{-}identification}\) then gives
\[
 \tau_Q^{\mathrm c}(x_0)=\tau_Q(x_0).                         \tag{27}
\]
Assumption \(\mathrm{ass{:}bounded\text{-}outcomes}\) gives \(-1\le Y(1)-Y(0)\le1\) almost surely, hence
\[
 -1\le\tau_Q^{\mathrm c}(x_0)\le1.                           \tag{28}
\]
On \(E_Q\), (27) places the causal target in the untruncated radius interval, while (28) places it in \([-1,1]\).  Definition \(\mathrm{def{:}private\text{-}confidence\text{-}interval}\) therefore gives
\[
 E_Q\subseteq\{\tau_Q^{\mathrm c}(x_0)\in C_n(h,k)\}.
\]
Taking probabilities and then the infimum over \(Q\in\mathcal P_{\mathrm{att}}\) yields
\[
 \inf_{Q\in\mathcal P_{\mathrm{att}}}
 \Pr_Q\{\tau_Q^{\mathrm c}(x_0)\in C_n(h,k)\}\ge1-\eta.
\]
The remainder lemma integrated its conditional pilot argument before use here, so this probability includes all three sample blocks, every Gaussian release, and the deterministic split.

%%% FIELD inclusion-statement
The concrete attainable model satisfies \(\mathcal P_{\mathrm{att}}\subset\mathcal P_{\mathrm{full}}\).

%%% FIELD inclusion-proof
Let \(Q\in\mathcal P_{\mathrm{att}}\).  By Definitions \(\mathrm{def{:}attainable\text{-}law\text{-}class}\) and \(\mathrm{def{:}full\text{-}law\text{-}class}\), every defining condition of \(\mathcal P_{\mathrm{full}}\), except possibly the two density conditions, is already satisfied by \(Q\).  Assumption \(\mathrm{ass{:}uniform\text{-}design}\) gives \(f_Q(x)=1\) for Lebesgue-almost every \(x\in\mathcal X\).  The declared parameter ranges \(0<\underline f\le1\) and \(1\le\overline f<\infty\) imply
\[
 f_Q(x)\le\overline f\quad\text{for Lebesgue-almost every }x\in\mathcal X
\]
and
\[
 f_Q(x)\ge\underline f\quad\text{for Lebesgue-almost every }
 x\in x_0+[-r_0,r_0]^d.
\]
Thus the global upper-density and anchor lower-density assumptions also hold.  Every member property in Definition \(\mathrm{def{:}full\text{-}law\text{-}class}\) is satisfied, so \(Q\in\mathcal P_{\mathrm{full}}\).  Since \(Q\) was arbitrary, \(\mathcal P_{\mathrm{att}}\subset\mathcal P_{\mathrm{full}}\).
